\chapter{תזמור סדרתי הלכה למעשה (Sequential Orchestration in Practice)}
\section{מדוע Sequential Process מתאים להפקת מסמכים}
הפקת מסמך עיון מקצועי אינה פעולה בודדת אלא רצף שלבים שבו כל שלב נשען על תוצר השלב שקדם לו, ולכן Sequential Process ב-CrewAI הוא הבחירה הטבעית עבורה. אי אפשר לחקור נושא לפני שקיים book plan הממקד את החקירה, אי אפשר לכתוב manuscript לפני שה-research pack מספק את החומר, אי אפשר לבצע review לפני שקיים draft, ואי אפשר לבנות latex spec לפני שהתוכן עבר ביקורת והתייצב. תלות חד-כיוונית זו היא בדיוק מה ש-Sequential Process מבטא: ה-Crew מריץ Task אחר Task בסדר קבוע, וה-Context של כל Task מורכב מתוצרי ה-Tasks שקדמו לו. היתרון המעשי הוא שהזרימה גלויה וניתנת לבדיקה, שכן בכל נקודה ברור איזה artifact כבר קיים ומהי הכניסה הצפויה לשלב הבא. בניגוד ל-hierarchical process, שבו Manager Agent מנתב עבודה באופן דינמי, ה-Sequential Process פשוט יותר להסבר, צפוי יותר בהתנהגותו, וקל יותר ל-observability ול-validation. בפרויקט זה הוא גם מבטא עיקרון הנדסי חשוב: עבודה דטרמיניסטית כמו citation management, יצירת assets, רינדור LaTeX והידור PDF נשארת מחוץ ל-Crew ורצה ב-Harness לאחר חמשת ה-Agent Tasks. כך נוצר גבול ברור בין שיקול דעת של מודל לבין צעדים ניתנים לאימות, וכל סטייה מהסדר הצפוי נחשפת מיד.

\section{חמשת ה-Agents המתמחים ותפקידיהם}
המערכת מורכבת מחמישה Agents, שכל אחד מהם מוגדר באמצעות role, goal ו-backstory, ובאופן מכוון תחומי האחריות שלהם אינם חופפים. ה-Planner Agent מתפקד כעורך טכני בכיר: הוא ממיר את הנושא ואת דרישות המטלה ל-book plan הכולל מבנה פרקים, מיקום של image, graph, table ו-formula, ו-acceptance checklist, אך אינו כותב פרוזה. ה-Research Agent הוא אנליסט מחקר האוסף מושגי מפתח, הגדרות ומועמדים למקורות, ורק לו מותר להשתמש בכלי חיפוש חיצוניים; הוא מציע source candidates אך אינו מייצר את ה-bibliography הסופי. ה-Writer Agent הוא כותב טכני בכיר העובד מתוך ה-Context בלבד, ללא חיפוש עצמאי, וממיר את התוכנית ואת המחקר ל-manuscript עם citation markers ו-placeholders. ה-Reviewer Agent הוא עורך ובוחן דרישות: הוא משפר בהירות ודיוק מבלי לשנות את המשמעות, ומפיק review report הכולל requirement checklist ואזהרות. ה-LaTeX Agent מתכנן את ה-assembly ומפיק latex spec, אך אינו מריץ קומפיילר ואינו אחראי ל-build correctness. הפרדה חדה זו, המעוגנת בסעיפי Not Responsible For עבור כל Agent, מונעת התנגשות תפקידים, שומרת על מערכת קטנה הניתנת להסבר, ומבטיחה שכל החלטה תיפול במקום אחד ויחיד בשרשרת. כך כל Agent נושא באחריות צרה וברורה, ושיתוף הפעולה ביניהם נשען על ממשק מוגדר ולא על חפיפה מקרית.

\section{העברת artifacts מובנים כ-Context לאורך השרשרת}
הדבק המחבר את חמשת ה-Tasks הוא ה-Context: כל Agent שלאחר ה-Planner צורך במפורש את תוצרי השלבים שקדמו לו, וההעברה נעשית באמצעות artifacts מובנים ולא כטקסט חופשי. ה-Planner מפיק את ה-book plan, המשמש Context ל-Research Task; ה-Research Agent מפיק את ה-research pack, וזה יחד עם ה-book plan מהווה את ה-Context של ה-Writing Task; ה-Writer מפיק את ה-manuscript; ה-Reviewer מקבל את שלושת התוצרים יחד ומחזיר reviewed manuscript ו-review report; ולבסוף ה-LaTeX Agent צורך את ה-book plan, את ה-reviewed manuscript ואת ה-review report כדי לבנות את ה-latex spec. שרשרת מצטברת זו מבטיחה שהחלטות מוקדמות, כמו ה-acceptance checklist של ה-Planner, נשארות זמינות ונבדקות עד סוף הזרימה. הבחירה ב-artifacts מובנים, רובם בפורמט JSON עם expected output מוגדר לכל Task, היא שמאפשרת validation דטרמיניסטי לאחר שה-Crew מסיים: ה-Harness יכול לוודא שכל citation marker מתאים למקור מוכר, שה-latex spec מפנה לקבצים תקפים, ושכל דרישת המטלה מכוסה. כך ה-Context אינו רק מנגנון תקשורת בין Agents אלא גם שכבת חוזה הניתנת לאכיפה, ההופכת את ה-Sequential Orchestration למערכת שקופה, ניתנת למעקב וניתנת לאימות מקצה לקצה.

לקריאה נוספת בנושאי פרק זה, ראו \cite{latex_project}.

